% GENERATED FILE. Edit canonical lessons, then run npm run generate:books.
% canonical-source-hash: fnv1a64:94b8c3ffd497610b
% canonical-lessons: GU-C82-seventeen, GU-C82-eighteen, GU-C82-twenty, GU-C82-thirty, GU-C82-forty

\chapter{Seventeen to Forty}
\label{ch:gu-seventeen-to-forty}
\hlchaptermodality{82}

\begin{chapteropening}
\textbf{By the end of this chapter:} \emph{I can count seventeen, eighteen, twenty, thirty and forty.}

Complete the last lesson of chapter 82: i can count seventeen, eighteen, twenty, thirty and forty.
\end{chapteropening}

\section[sattar]{\gu{સત્તર} (sattar) — seventeen}

\label{lesson:GU-C82-seventeen}



\begin{quote}
Before the new one: say the Gujarati for fifteen, then the Gujarati for sixteen.
\end{quote}

\subsection*{You'll want to know: \gu{સત્તર}}
\textbf{\gu{સત્તર}} (\emph{sattar}) — \textquotedblleft{}seventeen\textquotedblright{}.

\begin{grammarlens}[title={What you've built}]
One more word for seventeen to forty.
\end{grammarlens}

\subsection*{Your turn}
\begin{itemize}
\raggedright
  \item \emph{Say it:} \emph{sattar}
  \item \emph{Say it:} \emph{sattar}, once more
  \item \emph{Say it:} say \emph{sattar}
  \item \emph{Recall:} say the Gujarati for fifteen, then the Gujarati for sixteen, then say \emph{sattar} again
\end{itemize}

\begin{tcolorbox}[breakable,colback=teal!4,colframe=teal!35!black,title={Before you move on}]
What does \gu{સત્તર} mean? (\textquotedblleft{}Seventeen\textquotedblright{}.) Say it once more.
\end{tcolorbox}
\section[aḍhār]{\gu{અઢાર} (aḍhār) — eighteen}

\label{lesson:GU-C82-eighteen}



\begin{quote}
Before the new one: say the Gujarati for sixteen, then the Gujarati for seventeen.
\end{quote}

\subsection*{You'll want to know: \gu{અઢાર}}
\textbf{\gu{અઢાર}} (\emph{aḍhār}) — \textquotedblleft{}eighteen\textquotedblright{}.

\begin{grammarlens}[title={What you've built}]
One more word for seventeen to forty.
\end{grammarlens}

\subsection*{Your turn}
\begin{itemize}
\raggedright
  \item \emph{Say it:} \emph{aḍhār}
  \item \emph{Say it:} \emph{aḍhār}, once more
  \item \emph{Say it:} say \emph{aḍhār}
  \item \emph{Recall:} say the Gujarati for sixteen, then the Gujarati for seventeen, then say \emph{aḍhār} again
\end{itemize}

\begin{tcolorbox}[breakable,colback=teal!4,colframe=teal!35!black,title={Before you move on}]
What does \gu{અઢાર} mean? (\textquotedblleft{}Eighteen\textquotedblright{}.) Say it once more.
\end{tcolorbox}
\section[vīs]{\gu{વીસ} (vīs) — twenty}

\label{lesson:GU-C82-twenty}



\begin{quote}
Before the new one: say the Gujarati for seventeen, then the Gujarati for eighteen.
\end{quote}

\subsection*{You'll want to know: \gu{વીસ}}
\textbf{\gu{વીસ}} (\emph{vīs}) — \textquotedblleft{}twenty\textquotedblright{}.

\begin{grammarlens}[title={What you've built}]
One more word for seventeen to forty.
\end{grammarlens}

\subsection*{Your turn}
\begin{itemize}
\raggedright
  \item \emph{Say it:} \emph{vīs}
  \item \emph{Say it:} \emph{vīs}, once more
  \item \emph{Say it:} say \emph{vīs}
  \item \emph{Recall:} say the Gujarati for seventeen, then the Gujarati for eighteen, then say \emph{vīs} again
\end{itemize}

\begin{tcolorbox}[breakable,colback=teal!4,colframe=teal!35!black,title={Before you move on}]
What does \gu{વીસ} mean? (\textquotedblleft{}Twenty\textquotedblright{}.) Say it once more.
\end{tcolorbox}
\section[trīs]{\gu{ત્રીસ} (trīs) — thirty}

\label{lesson:GU-C82-thirty}



\begin{quote}
Before the new one: say the Gujarati for eighteen, then the Gujarati for twenty.
\end{quote}

\subsection*{You'll want to know: \gu{ત્રીસ}}
\textbf{\gu{ત્રીસ}} (\emph{trīs}) — \textquotedblleft{}thirty\textquotedblright{}.

\begin{grammarlens}[title={What you've built}]
One more word for seventeen to forty.
\end{grammarlens}

\subsection*{Your turn}
\begin{itemize}
\raggedright
  \item \emph{Say it:} \emph{trīs}
  \item \emph{Say it:} \emph{trīs}, once more
  \item \emph{Say it:} say \emph{trīs}
  \item \emph{Recall:} say the Gujarati for eighteen, then the Gujarati for twenty, then say \emph{trīs} again
\end{itemize}

\begin{tcolorbox}[breakable,colback=teal!4,colframe=teal!35!black,title={Before you move on}]
What does \gu{ત્રીસ} mean? (\textquotedblleft{}Thirty\textquotedblright{}.) Say it once more.
\end{tcolorbox}
\section[cāḷīs]{\gu{ચાળીસ} (cāḷīs) — forty}

\label{lesson:GU-C82-forty}



\begin{quote}
Before the new one: say the Gujarati for twenty, then the Gujarati for thirty.
\end{quote}

\subsection*{You'll want to know: \gu{ચાળીસ}}
\textbf{\gu{ચાળીસ}} (\emph{cāḷīs}) — \textquotedblleft{}forty\textquotedblright{}.

\begin{grammarlens}[title={What you've built}]
One more word for seventeen to forty.
\end{grammarlens}

\subsection*{Your turn}
\begin{itemize}
\raggedright
  \item \emph{Say it:} \emph{cāḷīs}
  \item \emph{Say it:} \emph{cāḷīs}, once more
  \item \emph{Say it:} say \emph{cāḷīs}
  \item \emph{Recall:} say the Gujarati for twenty, then the Gujarati for thirty, then say \emph{cāḷīs} again
\end{itemize}

\begin{tcolorbox}[breakable,colback=teal!4,colframe=teal!35!black,title={Before you move on}]
What does \gu{ચાળીસ} mean? (\textquotedblleft{}Forty\textquotedblright{}.) Say it once more.
\end{tcolorbox}
